\subsection{Same Cardinality, Finite, Countable, and Uncountable}

\begin{definition} \label{def:cardinality-equivalence} \hlt{Same Cardinality}\\
Two sets $A$ and $B$ have \hlt{same cardinality}, or are \hlt{equivalent}, if $\exists$ a bijection $f:A\to B$. Then $A\sim_{\#}B$.\\
The relation $\sim_{\#}$ is an equivalence relation: $A\sim_{\#}A$ by identity map, symmetry follows from Theorem \ref{thm:bijective-iff-invertible-function}, and transitivity follows by composing bijections.
\end{definition}

\begin{definition} \label{def:finite-countable-uncountable} \hlt{Finite, Countable, and Uncountable}\\
Set $J_0\defeq\emptyset$, and for $n\geq1$ set $J_n\defeq\{1,2,\ldots,n\}$. A set $A$ is
\begin{enumerate}[label=(\roman*)]
\setlength{\itemsep}{0pt}
\item \hlt{finite} if $A\sim_{\#} J_n$ for some $n\in\N$;
\item \hlt{infinite} if it is not finite;
\item \hlt{countable} if $A\sim_{\#} \N$;
\item \hlt{uncountable} if it is neither finite nor countable;
\item \hlt{at most countable} if it is finite or countable.
\end{enumerate}
Countable sets are also called enumerable or denumerable.
\end{definition}
